\subsection{Security und Konfigurationsgrenzen}
\label{subsec:security-konfigurationsgrenzen}

Tooling-Tests belegen die Konfigurationspriorität sowie die Validierung von
Umgebungen, Ports und CORS-Ursprüngen. Sie prüfen auch die profilabhängige
\path{.env.example} und \path{DATABASE_URL} nur für PostgreSQL. Lokal bestanden
diese Tests, während der externe Tooling-Job auf demselben Commit scheiterte
\parencite{kleiveist2026configuration,kleiveist2026ci-core-run}.

Nur \path{VITE_API_BASE_URL} darf in das Frontend gelangen. Tests verwerfen
öffentliche Secrets, entfernen \path{DATABASE_URL} aus Frontend und Tauri und
maskieren Datenbankpasswörter in Diagnosen. Ohne PostgreSQL entsteht keine
Datenbankvariable. Weitere Prüfungen decken Health, Readiness,
Fehlerantworten, datenbankunabhängige Liveness sowie Tauri-CSP und
Capability-Metadaten ab
\parencite{kleiveist2026acceptance,kleiveist2026release}.

Verifiziert sind damit nur Konfigurations-, Secret- und Desktop-Grenzen der
Baseline. Authentifizierung, Berechtigungen, produktive Credentials, TLS und
providerspezifische Härtung bleiben produktspezifisch; eine vollständige
Sicherheitsbewertung ist nicht belegt.
\beginnerinput{workspace/statement/600_qualitaet_verifikation/650}
